\documentclass[11pt,a4paper]{article}
\usepackage[utf8]{inputenc}
\usepackage[T1]{fontenc}
\usepackage[spanish]{babel}
\usepackage{hyperref}
\usepackage{geometry}
\geometry{margin=2.5cm}
\usepackage{parskip}

\title{330 el miedo es un costo, no una virtud}
\author{Alejandro Valenzuela Rivera\\\small Universidad Autónoma de Nuevo León}
\date{\today}
\begin{document}
\maketitle

\begin{quote}
\textbf{Sesgo declarado.} Este trabajo se sostiene desde el sesgo declarado de su autor, y no finge neutralidad. Lo que afirma se puede tumbar, y ese es el punto: el falsificador está escrito, no escondido.
\end{quote}

\section*{Resumen}
El miedo no se quita: lo que se queda es la cuenta. Sin miedo no hay especie que dure.

\section{Por qué se sostiene}
El que no tiene miedo no tiene costo: y una especie sin costo no mide lo que arriesga.

\section{De dónde viene}
'show me a species that is fearless and I'll show you something on the extinction list'.

\section{El caso}
Sinestro: la frase que dice que una especie sin miedo está en la lista de extinción.

\section{Método}
\textquestiondown{}El miedo se volvió cuenta, o se volvió virtud?

\section{Límites, y qué lo tumba}
Una especie sin miedo que sobreviva midiendo su costo.

\section{Lo que queda abierto}
\textquestiondown{}Se corrió contra la cuenta, o se suprimía el miedo?

\section*{Cómo se leyó esta tesis}
Esto se escribió en el repositorio man-made-empathy, del que este candado es una pieza: candados/06 del caso/330 el miedo es un costo, no una virtud.md. La máquina que lo compuso lo hizo desde este mismo archivo, y no agregó nada que no estuviera escrito. Si algo de esto es falso, se dice en el mismo sitio, y se corrige con fecha.

\end{document}
